\chapter{Literature Review}

\section{Introduction}

This chapter reviews existing academic literature on artificial intelligence (AI) and generative AI in marketing, with a particular focus on consumer perception. It aims to identify key themes and theoretical perspectives that explain how AI influences consumer attitudes and responses.

\section{Artificial Intelligence in Marketing}

Artificial intelligence has become an important tool in marketing. It is widely used in recommendation systems, personalised advertising, chatbots, and customer analytics. AI allows firms to process large amounts of data and deliver more targeted and efficient marketing strategies.

However, the use of AI also changes how consumers interact with marketing content. Instead of human-driven communication, many interactions are now mediated by algorithms and automated systems.

\section{Generative AI in Marketing}

Generative AI refers to technologies that can create content such as text, images, and videos. In marketing, it is used to generate advertisements, product descriptions, and social media content.

Compared to traditional AI, generative AI plays a more direct role in shaping marketing communication. This raises new questions about authenticity, creativity, and consumer trust.

\section{Consumer Perception}

Consumer perception refers to how individuals interpret and evaluate marketing information. It includes cognitive and emotional responses such as trust, perceived quality, credibility, and attractiveness.

In AI-driven marketing, perception becomes more complex because consumers may not always know whether content is created by humans or machines.

\section{Key Themes in AI Marketing and Consumer Perception}

Based on existing literature, several key themes can be identified:

\subsection{Trust}

Trust is one of the most important factors in AI marketing. Consumers may trust AI systems for their efficiency, but may also question their transparency and intentions.

\subsection{Authenticity}

AI-generated content can reduce perceived authenticity, especially when emotional or human-like communication is involved.

\subsection{Personalisation}

AI enables highly personalised marketing, which can improve relevance and satisfaction. However, excessive personalisation may also raise concerns.

\subsection{Privacy Concerns}

Consumers may feel uncomfortable when their data is used for AI-driven marketing. Privacy concerns can negatively affect attitudes and trust.

\subsection{Purchase Intention}

AI influences consumer decision-making, especially through recommendations and targeted advertising.

\section{Research Gap}

Although existing studies have explored AI in marketing, there are still gaps in understanding how generative AI specifically affects consumer perception. In particular, more research is needed to examine the interaction between trust, authenticity, and personalisation.

\section{Summary}

This chapter has reviewed key literature on AI and generative AI in marketing. It has identified major themes related to consumer perception, which will be further analysed in the methodology and findings chapters.